% Generated draft: requires mathematical and source review.
\documentclass[11pt]{article}
\usepackage[margin=25mm]{geometry}
\usepackage{fontspec}
\setmainfont{Latin Modern Roman}
\usepackage{amsmath,amssymb,mathtools,unicode-math}
\setmathfont{Latin Modern Math}
\usepackage{enumitem,needspace,xurl,hyperref}
\setlength{\parindent}{0pt}
\setlength{\parskip}{5pt}
\setlength{\emergencystretch}{3em}
\newcommand{\N}{\mathbb{N}}
\newcommand{\Z}{\mathbb{Z}}
\newcommand{\Q}{\mathbb{Q}}
\newcommand{\R}{\mathbb{R}}
\newcommand{\C}{\mathbb{C}}
\newcommand{\F}{\mathbb{F}}
\newcommand{\A}{\mathbb{A}}
\newcommand{\PP}{\mathbb{P}}
\newcommand{\E}{\mathbb{E}}
\newcommand{\eps}{\varepsilon}
\newcommand{\dd}{\,\mathrm{d}}
\begin{document}
\section*{UnsolvedMath ID: 30002331}
\textbf{Source problem code:} catalogue entry
\subsection{Definitions and mathematical statement}
The supplied record poses a research programme rather than a single conjecture. The following terminology makes a conventional interpretation explicit; the record itself does not specify the knot category or the precise conventions for its invariants.

In the conventional setting of tame knots in \(S^3\), represented by diagrams in the plane, an \(n\)-crossing diagram, for an integer \(n\geq 2\), is a knot diagram whose only singularities are transverse intersections of exactly \(n\) local branches, with pairwise distinct tangent lines and with the height ordering of the branches specified at each intersection. If \(\mathcal D_n(K)\) denotes the set of such diagrams representing a knot \(K\), define, whenever this set is nonempty,
\[
 c_n(K)=\min_{D\in\mathcal D_n(K)}\#\{\text{crossings of }D\}.
\]
Under this convention, the multi-crossing spectrum is the indexed collection \(\bigl(c_n(K)\bigr)_{n\geq 2}\), to the extent that its entries are defined. In particular, \(c_2(K)\) is the ordinary crossing number. The supplied record does not specify the index range, whether diagrams with no crossings are permitted, or how an empty admissible class is treated; these conventions require confirmation before this interpretation is adopted.

A petal projection conventionally means a knot projection with a single multiple crossing such that the arcs outside that crossing form loops, called petals, with no loop nested inside another. The crossing includes the height information needed to recover the knot. The record does not specify a procedure that turns such a projection into a braid, nor an equivalence relation under which the resulting braids are to be classified.

\textbf{Research questions.} For knots \(K\) in the intended category, describe the behavior of their multi-crossing spectra; identify the knot families that arise naturally from these spectra; describe the braids arising from petal projections under the intended construction; and determine relationships between the resulting invariants and traditional knot invariants.

The phrases ``behavior,'' ``arise naturally,'' and ``traditional knot invariants'' have no specified mathematical targets in the supplied statement. No particular inequality, characterization, family, braid construction, or list of comparison invariants is supplied. Consequently, the record does not determine a unique quantified theorem or classification problem. The definitions above provide a provisional framework, not additional hypotheses or conjectures attributed to the source.
\subsection{Short English statement}
Study multi-crossing spectra of knots, the knot families associated with them, braids obtained from petal projections, and connections with traditional knot invariants. The supplied question is a broad research programme: its diagram conventions, braid construction, equivalence relation, and precise classification targets require clarification. The supplied status is \emph{partially solved}, according to dated curation rather than independent verification.
\subsection{Sources}
\begin{itemize}
\item Local UnsolvedMath registry: supplied complete source record for permanent ID \(30002331\), title \emph{Spectra of Multi-Crossing Knot Numbers}, exact source problem code \texttt{catalogue entry}, in the \emph{Oberwolfach Reports Open Problems} collection. The proposed year is recorded as 2013; no individual proposer is identified. The registry labels the statement \texttt{exact} and reports a prior manual review for self-containment, coherence, and extraction artifacts. These are supplied registry annotations, not verification performed for this draft.
\item Supplied report citation: \emph{Geometric Knot Theory (2014)}, Oberwolfach Reports, pp.\ 1313--1358, DOI \texttt{10.4171/owr/2013/22}. \url{https://doi.org/10.4171/owr/2013/22}. The literature list also calls this source \emph{Oberwolfach Report 22/2013}. The supplied 2014 title date and 2013 report designation are retained without reconciliation.
\item Supplied literature reference: Adams et al., \emph{Multi-crossing Number for Knots and the Kauffman Bracket Polynomial}. \url{https://doi.org/10.1017/S0305004116000906}.
\item Literature assessment dated \texttt{2026-08-21}: the supplied curation labels the problem \emph{partially solved} and reports monotonicity, composition results, and extensive computations for multi-crossing spectra, while describing the general spectrum behavior and the resulting braid classification as remaining open. This assessment is supplied dated curation derived from the OWR record, not independently verified. Its time-sensitive status and cited claims require independent checking.
\end{itemize}
\end{document}
